\section{Results}\label{sec:results}

\subsection{Main estimates}\label{sec:results-main}

\Cref{tab:main} reports the estimates. Conditional on $x$, treatment raises the
outcome by \val{beta_treated} (standard error \val{se_treated}), which is
\val{beta_treated_pct_mean} percent of the control mean.

\begin{table}[htbp]
\centering
\begin{threeparttable}
\caption{Effect of treatment on the outcome}\label{tab:main}
\exhibit{tables/tab_example.tex}
\begin{tablenotes}\small
\item \textit{Notes:} OLS estimates; standard errors in parentheses. Sample of
\val{n_units} simulated units. Source: \path{code/analysis/tab_example.R}.
\end{tablenotes}
\end{threeparttable}
\end{table}

\subsection{Heterogeneity}\label{sec:results-heterogeneity}

\Cref{fig:main} plots the outcome against the covariate by treatment status.

\begin{figure}[htbp]
\centering
\exhibitgraphic[width=0.8\textwidth]{figures/fig_example.pdf}
\caption{Outcome and covariate by treatment status}\label{fig:main}
\begin{minipage}{0.9\textwidth}\footnotesize
\textit{Notes:} Each point is a unit; lines are OLS fits within group.
Source: \path{code/analysis/fig_example.R}.
\end{minipage}
\end{figure}

% An explicit passage is a finer unit than a subsection: it is tracked, and
% signed off, on its own.
\begin{passage}{results-magnitude}[code/analysis/tab_example.R]
The slope is parallel across groups, so a single treatment coefficient of
\val{beta_treated} summarises the shift.
\end{passage}
